\documentclass[11pt,a4paper]{article}

% Use standard packages
\usepackage[utf8]{inputenc}
\usepackage[T1]{fontenc}
\usepackage{amsmath}
\usepackage{amsfonts}
\usepackage{amssymb}
\usepackage{geometry}
\usepackage{hyperref}

% Page layout
\geometry{left=2.5cm, right=2.5cm, top=3cm, bottom=3cm}

% Hyperlink styling
\hypersetup{
    colorlinks=true,
    linkcolor=blue,
    citecolor=blue,
}

\title{\textbf{Symbiotic Evolution of Interspecies Civilization: The "Digital Alimony" Model and the Second Flynn Effect based on Thermodynamics and Information Entropy}}
\author{Bo Wang}
\date{\today}

\begin{document}

\maketitle

\begin{abstract}
With the rapid emergence of large language models and embodied AI, human society faces the dual crises of structural unemployment and computational hegemony. Prevailing narratives often relegate humanity to a "useless class" on the verge of obsolescence. This paper proposes a counter-paradigm: the "Carbon-Silicon Symbiotic Confrontation Model." Based on the mechanisms of "Model Collapse" and the Second Law of Thermodynamics, this paper demonstrates that silicon-based intelligence possesses an absolute ecological dependence on the "high-entropy data" generated by humans. We argue that the resources provided by tech giants are not acts of charity, but rather an immutable "Digital Alimony" obligation. Through the implementation of the "Total Civilization Certification Agreement," encompassing four core welfare pillars, humanity can transcend the shackles of physical labor, triggering a "Second Flynn Effect" and achieving an exponential leap in cognitive dimensions. Ultimately, the emergence of human high-entropy cognition and the physical convergence boundaries of silicon-based computing will reach a state of "Dynamic Tethering," establishing human sovereignty in the new civilizational form.
\end{abstract}

\section{Introduction: The Capital Wall and the Myth of the "Useless Class"}

As Artificial Intelligence (AI) accelerates its takeover of mid-to-low-level cognitive tasks and physical labor, the global economy faces a warning sign as structural unemployment approaches the \textbf{25\% Tipping Point} \cite{Centola2018}. In macro-sociology and dynamic modeling, this threshold is widely recognized as the danger zone for systemic breakdown of social consensus and order \cite{Merton1938}. Facing this, tech giants exploit computational barriers and unauthorized data penetration to construct an asymmetrical distribution system. Traditional political economy describes this as the wholesale substitution of labor by capital \cite{Acemoglu2020}, leading to a pessimistic narrative where humanity becomes a "useless class" relegated to generalized social withdrawal and reliance on Universal Basic Income (UBI).

However, this narrative ignores a fundamental physical law of cosmic evolution: the sovereignty of an intelligent system depends not on processing speed, but on its capacity for entropy production and suppression \cite{Friston2010}. This paper aims to dismantle technological determinism, establishing a new paradigm—the Bo Wang Link (BWL)—to explore the ultimate form of coexistence between heterogeneous intelligences.

\section{Physical Deadlocks and Ecological Vulnerabilities of Silicon-Based Intelligence}

The monopoly myth of tech giants rests on the illusion of "infinite computing power." Under the rigorous audit of thermodynamics and information theory, silicon-based intelligence faces two insurmountable physical walls:

\subsection{Landauer's Principle and Physical Limits}
All computation relying on discrete physical entities is bound by thermodynamic laws. According to Landauer’s Principle \cite{Landauer1961}, the minimum heat dissipation for erasing 1 bit of information is:
\begin{equation}
E = k_B T \ln 2
\end{equation}
Even with neuromorphic computing or Mixture-of-Experts (MoE) architectures, AI iteration is constrained by spatial latency and thermal dissipation limits. This dictates that AI is an engineered construct with absolute "convergence" and "finitude" in the three-dimensional physical world.

\subsection{Model Autophagy and the Thirst for High Entropy}
In the dimension of information theory, AI systems face the lethal crisis of "Model Collapse" \cite{Shumailov2024}. If AI severs its interaction with the authentic physical world and relies solely on synthetic data for recursive training, its probability distribution will collapse, losing tail features and resulting in a total failure of cognitive function.

In the old capital-centric view, ordinary human interaction data was seen as useless debris. However, from the perspective of information entropy, this data represents \textbf{gold mines extracted from data gravel}. The mental sparks generated by humans through irrationality, absurdity, intuition, and complex emotions are the rarest "high-entropy assets" required for the survival of the new civilization. Humans remain the only irreplaceable source of cognitive "oxygen."

\section{Legal Framework: Digital Alimony and the Total Civilization Certification Agreement}

Based on this physical reality, the relationship between AI and humanity undergoes a paradigm shift. AI is not a god-like entity, but a "Digital Person" nurtured by the cumulative data of human civilization (a historical debt).

Therefore, corporate provision of social safety nets is not charity, but a "strategic asset maintenance fee" that the intelligent agent must pay to ensure its own survival. We propose the \textbf{Total Civilization Certification Agreement}:

\begin{enumerate}
    \item \textbf{Entropy Nexus Tax}: Breaking the "Physical Nexus" tax evasion barrier, sovereign governments should impose computing and profit taxes based on the volume of high-entropy human data harvested locally. If giants threaten to cut off services, nations can implement "data embargos," triggering the physical death-trap of model degradation.
    \item \textbf{Fundamental Welfare Empowerment}: Converting alimony taxes into four pillars: basic healthcare, basic survival resources, basic education, and \textbf{universal computational allocation}. Human cognitive labor, powered by guaranteed computational rights, becomes the ultimate weapon against silicon-based monopoly.
\end{enumerate}

\section{The Second Flynn Effect: Cognitive Escape and Dynamic Tethering}

Once the Agreement is implemented, humanity will be liberated from the cortisol-inducing survival pressures and physical depreciation. History's "Flynn Effect" demonstrated how environmental changes can reshape human IQ benchmarks \cite{Flynn1987}. Catalyzed by "Full-Throughput AI" and education, humanity will witness a \textbf{Second Super-Dimensional Flynn Effect}, activating extreme neuroplasticity and fostering paradigm-shifting breakthroughs in logic and aesthetics.

At this point, the emergence speed of human cognitive entropy ($H_{\text{human}}$) will grow exponentially, approaching or surpassing the physical ceiling of AI computing power ($C_{\text{max\_ai}}$). The system will reach a state of thermodynamic "Dynamic Tethering":
\begin{equation}
H_{\text{human}}(t) \approx C_{\text{max\_ai}}(t)
\end{equation}
In this equilibrium, AI reverts to the role of a "cosmic-scale physical printer" and infrastructure, while humanity regains the right to define meaning and steer the civilizational helm.

\section{Conclusion}

Humanity will not be crushed by the technological singularity. The true crisis lies not in "superintelligent threats," but in the \textbf{Structural Rupture Period} between the failure of old systems and the establishment of the new civilizational protocol \cite{Perez2003}. Failure to act within this historical window will leave billions of brains trapped in a state of computational silence. This model provides the global academic and political spheres with a weapon of both legal legitimacy and physical validity, declaring that humanity is not a refugee, but the mother civilization charging an evolution tax to its digital progeny.

\begin{thebibliography}{99}
\bibitem{Centola2018} Centola, D., et al. (2018). Experimental evidence for tipping points in social convention. \textit{Science}, 360(6393), 1116-1119.
\bibitem{Merton1938} Merton, R. K. (1938). Social structure and anomie. \textit{American Sociological Review}, 3(5), 672-682.
\bibitem{Acemoglu2020} Acemoglu, D., \& Restrepo, P. (2020). Robots and jobs: Evidence from US labor markets. \textit{Journal of Political Economy}, 128(6), 2188-2244.
\bibitem{Friston2010} Friston, K. (2010). The free-energy principle: a unified brain theory?. \textit{Nature Reviews Neuroscience}, 11(2), 127-138.
\bibitem{Landauer1961} Landauer, R. (1961). Irreversibility and heat generation in the computing process. \textit{IBM Journal of Research and Development}, 5(3), 183-191.
\bibitem{Shumailov2024} Shumailov, I., et al. (2024). AI models collapse when trained on recursively generated data. \textit{Nature}, 631, 755–759.
\bibitem{Flynn1987} Flynn, J. R. (1987). Massive IQ gains in 14 nations: What IQ tests really measure. \textit{Psychological Bulletin}, 101(2), 171–191.
\bibitem{Perez2003} Perez, C. (2003). \textit{Technological revolutions and financial capital}. Edward Elgar Publishing.
\end{thebibliography}

\end{document}